\addchap{A Anmerkung zur Nutzung von Künstlicher Intelligenz}
\setcounter{chapter}{1}

Im Rahmen dieser Arbeit wurden Künstliche Intelligenz (\acrshort{acr:KI})\index{Künstliche Intelligenz} basierte Werkzeuge benutzt. Der folgende Absatz gibt eine Übersicht über die verwendeten Werkzeuge und den jeweiligen Einsatzzweck.

Im Software / Elektronik Bereich wurde die KI Gemini zur Coding-Hilfe und Recherche benutzt. Bei der Erstellung dieser Dokumentation wurde ebenfalls Gemini benutzt. Die Nutzung bestand vor allem darin die Sprache in allen Bereichen an den wissenschaftlichen Stil anzupassen, wo dies noch nicht geschehen war. Vor allem im Bereich Software-Dokumentation und in Bezug auf das Projektmanagement (Tracebility, Testing) kam die KI auch zur Generierung von Inhalt zum Einsatz. Speziell wurden auch alle UML-Digramma mithilfe von Plant-UML mit KI erstellt.

\addchap{B Ergänzung}
\setcounter{chapter}{2}
Was noch wichtig zu erwähnen ist, ist dass bei der Entwicklung für den STM32-Mikrocomputer neben der Cube-IDE auch CubeMX verwendet wurde. Das erleichterte die Implementierung von Funktionen wie PWM und die Interrupts.
